\section{The household counterpart: ownership without access}
\label{sec:household}

The revenue side of the condition fails because of who owns the corporate claims. This
section asks whether the same ownership structure governs the household side, measured on
survey microdata rather than on tax components. It does, and the mechanism is the same.

\subsection{Whose capacity falls}

In the Survey of Consumer Finances we move a share of aggregate wage income to capital
income, distribute it by actual ownership and hold debts fixed: incidence under stated
conventions, not a forecast. With payments fixed, a household's burden rises exactly when
its income falls.

Three results hold across the conventions we varied. Capacity falls for about \result{HHBottomQuartileAffected} percent
of bottom-wealth-quartile households against about \result{HHTopOnePctAffected} percent of
the top one percent, monotone in wealth. About three fifths of household debt is owed by
households earning more from wages than they hold in capital, at any size of shift. And
broadened ownership through retirement vehicles moves nothing unless the holdings can be
spent, because ownership without access cannot service a payment:
\result{HHEquityOutsideHouseholds} percent of corporate equity sits outside the household
sector, another \result{HHEquityInRetirement} percent in retirement vehicles.

The fiscal condition falls short on a rate, which a legislature can change; this one on who
holds the capital, which a rate does not. Closing the whole fiscal gap and distributing it
per capita moves household capacity little.

Illustrative only, and post hoc: for the median dollar of affected debt, restoring capacity
takes roughly \result{HHRestoreFive} percent cumulative real growth at a five percent shift,
\result{HHRestoreTen} percent at ten. Median and upper quartile are stable across the 2019
and 2022 waves; the lower quartile is not. The pre-specified boundary is a maximum set by
one household and is not reported. Nor are the policy arrangements, since one statistic
counts households crossing a line and the other conditions on a group the policy changes,
nor the holder mapping, which moves sevenfold with the agency classification. It is neither
the debt-service stress test of \citet{BartscherKuhnSchularickSteins2020} nor the
model-based growth threshold of \citet{Bayraktar2026}.

\subsection{Exposure type is a pay effect, and how displacement arrives changes who pays}

Households whose earners are physically exposed look more fragile than those whose earners are
cognitively exposed, and that difference does not survive a control for pay: reweighting either
group to the other's pay distribution leaves \result{PayGapSurvivingLow} to
\result{PayGapSurvivingHigh} percent of the raw gap, and on one index and one survey the contrast
reverses sign. Exposure type is largely a proxy for pay here, and the direction, confirmed by the
rebuild, is what we report.


